% GENERATED FILE. Edit canonical lessons, then run npm run generate:books.
% canonical-source-hash: fnv1a64:dda408c0fd1c61ab
% canonical-lessons: GE-C126-brieftraeger, GE-C126-kellner, GE-C126-kunde, GE-C126-kasse, GE-C126-chef

\chapter{Postman, Waiter, Customer, Till, Boss}
\label{ch:ge-postman-waiter-customer-till-boss}
\hlchaptermodality{126}

\begin{chapteropening}
\textbf{By the end of this chapter:} \emph{I can say a postman, a waiter, a customer, a till and a boss.}

Complete the last lesson of chapter 126: I can say a postman, a waiter, a customer, a till and a boss.
\end{chapteropening}

\section[der Briefträger]{der Briefträger — a postman}

\label{lesson:GE-C126-brieftraeger}



\begin{quote}
Before the new one: say the German for a lawyer, then the German for an engineer.
\end{quote}

\subsection*{You'll want to know: der Briefträger}
\textbf{der Briefträger} — \textquotedblleft{}a postman\textquotedblright{}.

\begin{cousinweb}
\textbf{Briefträger} is the word for \textquotedblleft{}letter\textquotedblright{} + \emph{Träger}, \textquotedblleft{}carrier\textquotedblright{}, from \emph{tragen}. The letter word $\leftarrow$ Old High German \emph{briaf} $\leftarrow$ Latin \emph{brevis}, \textquotedblleft{}short (document)\textquotedblright{} (EtymDB 2.1).
\end{cousinweb}

\begin{grammarlens}[title={What you've built}]
One more word for this chapter.
\end{grammarlens}

\subsection*{Your turn}
\begin{itemize}
\raggedright
  \item \emph{Say it:} \emph{der Briefträger}
  \item \emph{Say it:} \emph{der Briefträger}, once more
  \item \emph{Say it:} say \emph{der Briefträger}
  \item \emph{Recall:} say the German for a lawyer, then the German for an engineer, then say \emph{der Briefträger} again
\end{itemize}

\begin{tcolorbox}[breakable,colback=teal!4,colframe=teal!35!black,title={Before you move on}]
What does der Briefträger mean? (\textquotedblleft{}A postman\textquotedblright{}.) Say it once more.
\end{tcolorbox}
\section[der Kellner]{der Kellner — a waiter}

\label{lesson:GE-C126-kellner}



\begin{quote}
Before the new one: say the German for an engineer, then the German for a postman.
\end{quote}

\subsection*{You'll want to know: der Kellner}
\textbf{der Kellner} — \textquotedblleft{}a waiter\textquotedblright{}.

\begin{cousinweb}
\textbf{Kellner} $\leftarrow$ Old High German \emph{kelnere} $\leftarrow$ Medieval Latin \emph{cellenārius}, \textquotedblleft{}cellar master\textquotedblright{} (EtymDB 2.1) — borrowed: the man who kept the wine cellar.
\end{cousinweb}

\begin{grammarlens}[title={What you've built}]
One more word for this chapter.
\end{grammarlens}

\subsection*{Your turn}
\begin{itemize}
\raggedright
  \item \emph{Say it:} \emph{der Kellner}
  \item \emph{Say it:} \emph{der Kellner}, once more
  \item \emph{Say it:} say \emph{der Kellner}
  \item \emph{Recall:} say the German for an engineer, then the German for a postman, then say \emph{der Kellner} again
\end{itemize}

\begin{tcolorbox}[breakable,colback=teal!4,colframe=teal!35!black,title={Before you move on}]
What does der Kellner mean? (\textquotedblleft{}A waiter\textquotedblright{}.) Say it once more.
\end{tcolorbox}
\section[der Kunde]{der Kunde — a customer}

\label{lesson:GE-C126-kunde}



\begin{quote}
Before the new one: say the German for a postman, then the German for a waiter.
\end{quote}

\subsection*{You'll want to know: der Kunde}
\textbf{der Kunde} — \textquotedblleft{}a customer\textquotedblright{}.

\begin{cousinweb}
\textbf{Kunde} $\leftarrow$ Old High German \emph{cundo}, \textquotedblleft{}one who is known\textquotedblright{}, from \emph{kund}, \textquotedblleft{}known\textquotedblright{} $\leftarrow$ Proto-Germanic \emph{*kunþaz} (EtymDB 2.1) — inherited. English \emph{uncouth}, \textquotedblleft{}unknown, strange\textquotedblright{}, shares the root.
\end{cousinweb}

\begin{grammarlens}[title={What you've built}]
One more word for this chapter.
\end{grammarlens}

\subsection*{Your turn}
\begin{itemize}
\raggedright
  \item \emph{Say it:} \emph{der Kunde}
  \item \emph{Say it:} \emph{der Kunde}, once more
  \item \emph{Say it:} say \emph{der Kunde}
  \item \emph{Recall:} say the German for a postman, then the German for a waiter, then say \emph{der Kunde} again
\end{itemize}

\begin{tcolorbox}[breakable,colback=teal!4,colframe=teal!35!black,title={Before you move on}]
What does der Kunde mean? (\textquotedblleft{}A customer\textquotedblright{}.) Say it once more.
\end{tcolorbox}
\section[die Kasse]{die Kasse — a till, a checkout}

\label{lesson:GE-C126-kasse}



\begin{quote}
Before the new one: say the German for a waiter, then the German for a customer.
\end{quote}

\subsection*{You'll want to know: die Kasse}
\textbf{die Kasse} — \textquotedblleft{}a till, a checkout\textquotedblright{}.

\begin{cousinweb}
\textbf{Kasse} $\leftarrow$ Italian \emph{cassa}, \textquotedblleft{}box\textquotedblright{} $\leftarrow$ Latin \emph{capsa}, \textquotedblleft{}case, box\textquotedblright{} (EtymDB 2.1) — borrowed. English \emph{cash} and \emph{case} come from the same word.
\end{cousinweb}

\begin{grammarlens}[title={What you've built}]
One more word for this chapter.
\end{grammarlens}

\subsection*{Your turn}
\begin{itemize}
\raggedright
  \item \emph{Say it:} \emph{die Kasse}
  \item \emph{Say it:} \emph{die Kasse}, once more
  \item \emph{Say it:} say \emph{die Kasse}
  \item \emph{Recall:} say the German for a waiter, then the German for a customer, then say \emph{die Kasse} again
\end{itemize}

\begin{tcolorbox}[breakable,colback=teal!4,colframe=teal!35!black,title={Before you move on}]
What does die Kasse mean? (\textquotedblleft{}A till, a checkout\textquotedblright{}.) Say it once more.
\end{tcolorbox}
\section[der Chef]{der Chef — a boss}

\label{lesson:GE-C126-chef}



\begin{quote}
Before the new one: say the German for a customer, then the German for a till.
\end{quote}

\subsection*{You'll want to know: der Chef}
\textbf{der Chef} — \textquotedblleft{}a boss\textquotedblright{}.

\begin{cousinweb}
\textbf{Chef} $\leftarrow$ French \emph{chef}, \textquotedblleft{}head\textquotedblright{} $\leftarrow$ Latin \emph{caput} (EtymDB 2.1) — borrowed. English \emph{chief} is the same word; English \emph{chef} kept only the kitchen sense.
\end{cousinweb}

\begin{grammarlens}[title={What you've built}]
One more word for this chapter.
\end{grammarlens}

\subsection*{Your turn}
\begin{itemize}
\raggedright
  \item \emph{Say it:} \emph{der Chef}
  \item \emph{Say it:} \emph{der Chef}, once more
  \item \emph{Say it:} say \emph{der Chef}
  \item \emph{Recall:} say the German for a customer, then the German for a till, then say \emph{der Chef} again
\end{itemize}

\begin{tcolorbox}[breakable,colback=teal!4,colframe=teal!35!black,title={Before you move on}]
What does der Chef mean? (\textquotedblleft{}A boss\textquotedblright{}.) Say it once more.
\end{tcolorbox}
